% ============================================================================
% 4.1 Verilog 基本概念
% ============================================================================
\subsection{软件和硬件的根本区别}

汇编语言是给 CPU 看的——它描述的是"一步一步做什么"。
Verilog 是给综合器看的——它描述的是"硬件模块之间如何连接、每个时钟沿做什么"。

\paragraph{并行 vs 顺序}
在 C/汇编里，代码是顺序执行的。Verilog 不一样：
每个 \texttt{always} 块都在\textbf{同时}运行。CPU 内部有几十上百个
这样的并行块。

\paragraph{两个基本类型}
\begin{itemize}
    \item \textbf{reg}：寄存器，在时钟沿保存值
    \item \textbf{wire}：连线，只是把一个模块的输出连到另一个模块的输入
\end{itemize}

\subsection{同步逻辑：时钟驱动一切}

几乎所有现代数字设计都使用"同步时序逻辑"：

\begin{lstlisting}[style=verilog]
always @(posedge clk or negedge rst_n)
begin
    if (!rst_n)
        q <= 1'b0;        // 复位
    else
        q <= d;           // 每个时钟上升沿，把 d 的值写入 q
end
\end{lstlisting}

\begin{itemize}
    \item \texttt{posedge clk}：时钟上升沿
    \item \texttt{negedge rst\_n}：复位（低电平有效）
    \item \texttt{<=}：非阻塞赋值，在组合逻辑里用 \texttt{=}，在时序逻辑里用 \texttt{<=}
    \item \texttt{1'b0}：表示 1 位二进制 0
\end{itemize}

\subsection{UART 通信的基本约定}

\begin{center}
\begin{tabular}{l l}
波特率 & 每秒发送的比特数，常见值如 115200 \\
数据位 & 一个字节通常 8 位 \\
停止位 & 1 位或 2 位，标识字节结束 \\
起始位 & 1 位，先把电平拉低来通知"我要开始发了" \\
\hline
总位数 & 1 起始 + 8 数据 + 1 停止 = 10 位 / 字节 \\
\end{tabular}
\end{center}

UART 不附带时钟信号，收发双方必须事先约定好同一个波特率，
然后自己在内部用时钟计数来估计每个比特的中心位置。
